\documentclass[12pt,oneside]{uhmasterthesis}

\usepackage[spanish,es-tabla]{babel}
\usepackage[utf8]{inputenc}
\usepackage[T1]{fontenc}
\usepackage{fvextra}
\usepackage{csquotes}
\usepackage{amsmath,amssymb,amsfonts,amsthm,mathtools}
\usepackage{graphicx}
\usepackage{array}
\usepackage{booktabs}
\usepackage{longtable}
\usepackage{microtype}
\usepackage[hypertexnames=false,hidelinks,linktocpage]{hyperref}
\usepackage{tikz}
\usetikzlibrary{positioning,fit,backgrounds,arrows.meta,calc}
\setlength{\emergencystretch}{2em}

\theoremstyle{definition}
\newtheorem{definition}{Definición}[chapter]
\newtheorem{property}{Propiedad}[chapter]
\newtheorem{auxrule}{Regla}
\renewcommand{\theauxrule}{A\arabic{auxrule}}

\title{Unified Multi-Agent Coordination Bridging Large Language Models and Symbolic AI for Autonomous Systems}
\author{Lic. Roberto Marti Cedeno}
\university{Universidad de La Habana}
\faculty{Facultad de Matemática y Computación}
\advisor{}
\coadvisor{}
\degree{Máster en Ciencia de la Computación}
\worktype{Tesis de Maestría}
\date{La Habana, junio de 2026}
\logo{Graphics/uhlogo}
\qrlogo{}
\repositoryurl{https://github.com/rmarticedeno/unified-multi-agent-coordination}

\hypersetup{
  pdftitle={Unified Multi-Agent Coordination Bridging Large Language Models and Symbolic AI for Autonomous Systems},
  pdfauthor={Lic. Roberto Marti Cedeno},
  pdfsubject={Tesis de Maestría en Ciencia de la Computación},
  pdfkeywords={sistemas multiagente, modelos grandes de lenguaje, inteligencia artificial simbólica, sistemas autónomos}
}

\begin{document}

\frontmatter
\maketitle
\chapter*{Abstract}
\addcontentsline{toc}{chapter}{Abstract}

Large language models (LLMs) enable flexible interpretation and task decomposition, but linguistic plausibility does not establish that a multi-agent plan is executable, authorized, or complete. Conversely, symbolic artificial intelligence provides explicit constraints and verifiable state transitions but is less suited to interpreting underspecified natural-language goals. This thesis addresses the problem of coordinating heterogeneous autonomous agents without allowing an LLM proposal to substitute for an executable justification. Following a design-science methodology, it develops a protocol-independent framework that separates linguistic proposal from symbolic authorization. The framework represents agent capabilities, task dependencies, authority constraints, commitments, artifact contracts, and completion evidence explicitly. It returns either an authorized coordination plan or a structured infeasibility report. The framework is instantiated as a Python \texttt{CoordinationAgent} over \texttt{CoordinationSdk} middleware: the coordination agent owns planning and authorization, while the SDK owns communication, runtime adaptation, and normalization of remote, local, and optional Lingo-based linguistic agents. Evaluation uses unit and integration tests, seven deterministic coordination scenarios, an eleven-check Docker A2A system harness, and batched local LLM reference runs with \texttt{qwen/qwen3-1.7b} and \texttt{google/gemma-4-e2b}. The evidence supports the main boundary claim: language can propose structured candidates, but symbolic checks authorize or refuse execution and all dispatch proceeds through the SDK-managed communication surface. The work does not claim broad benchmark superiority over all possible baselines; its principal contribution is a precise and reproducible operational boundary for hybrid LLM--symbolic coordination.

\tableofcontents
\cleardoublepage
\listoffigures
\cleardoublepage
\listoftables

\mainmatter
\chapter{Introduction}

Autonomous systems increasingly require coordination across heterogeneous agents, uncertain operating conditions, and mixed reasoning demands. Large language models provide flexible natural language understanding and planning assistance, while symbolic AI contributes explicit structure, verifiability, and rule-based reasoning.

This chapter introduces the problem context, research motivation, objectives, and thesis organization.

\chapter{Background}

This chapter reviews the foundations relevant to unified multi-agent coordination, including multi-agent systems, large language models, symbolic AI, hybrid reasoning architectures, and autonomous system design.

The final thesis should use this chapter to position the proposed work relative to prior research and current implementation practices.

\chapter{Metodología de investigación}
\label{chap:methodology}

\section{Paradigma de investigación y tipo de estudio}

This thesis adopts a design-science paradigm because its central research output is an artifact: a coordination framework with explicit constructs, decision rules, an architectural instantiation, and an evaluation method. The study combines exploratory and analytical work in the literature review, formal and architectural design in the construction phase, and quantitative and qualitative analysis in the planned evaluation. Quantitative measures assess classification, task execution, latency, and overhead; qualitative trace analysis examines whether authorization and refusal decisions are intelligible and supported by recorded evidence.

\section{Proceso de investigación}

The research process contains five linked stages. First, a critical literature review identifies the strengths and unresolved limitations of classical multi-agent coordination, symbolic reasoning, LLM agents, and interoperability protocols. Second, the research gap is translated into requirements for openness, interoperability, bounded authority, traceability, and explicit refusal. Third, these requirements are formalized as a protocol-independent framework. Fourth, the framework is mapped to a Python \texttt{CoordinationAgent} design over SDK middleware that uses A2A as its interoperability backbone. Fifth, controlled scenarios compare the complete design with baselines and ablations. Each stage produces evidence for one or more specific objectives rather than serving as an independent implementation activity.

\begin{table}[htbp]
\centering
\caption{Alineación de objetivos, métodos y evidencia.}
\label{tab:method-alignment}
\begin{tabular}{p{0.12\textwidth}p{0.38\textwidth}p{0.38\textwidth}}
\toprule
\textbf{Objetivo} & \textbf{Método} & \textbf{Evidencia} \\
\midrule
1 & Síntesis temática crítica de la literatura & Criterios de comparación y brecha explícita de investigación \\
2 & Modelado formal y conceptual & Definiciones, reglas, predicados de viabilidad y propiedades declaradas \\
3 & Diseño arquitectónico y mapeo de protocolo & Componentes, registros, flujos de trabajo y trazabilidad teoría-implementación \\
4 & Evaluación controlada por escenarios & Comparaciones con líneas base, ablaciones, métricas y trazas de ejecución \\
5 & Discusión analítica & Análisis de alcance, casos de fallo, limitaciones y amenazas a la validez \\
\bottomrule
\end{tabular}
\end{table}

\section{Unidad de análisis y diseño de escenarios}

The primary unit of analysis is a coordination attempt: one problem request, one admitted registry snapshot, the candidate plan or refusal produced for that request, and the resulting execution trace when a plan is authorized. Scenarios are constructed by varying four factors: capability coverage, artifact compatibility, authority constraints, and runtime availability. This factorial view prevents evaluation from relying only on successful demonstrations.

The scenario suite must contain feasible direct-delegation cases, feasible decomposed cases, cases requiring an approved bounded linguistic transformation, and infeasible cases caused by missing capability, incompatible modality, cyclic dependency, insufficient authority, or unavailable agents. Each scenario includes a reference feasibility label and the minimum justification required for that label. Reference labels should be reviewed independently before system outputs are examined.

\section{Líneas base y ablaciones}

Three configurations provide the minimum comparison. The \emph{hybrid CoordinationAgent} uses \texttt{CoordinationSdk} as the only communication boundary, may invoke Lingo-backed SDK capabilities for structured interpretation and candidate decomposition, and always applies symbolic feasibility checks before dispatch. The \emph{CoordinationAgent without symbolic feasibility} variant receives the same request and registry descriptions but accepts or rejects its own proposed plan without the deterministic feasibility layer. The \emph{rule-only CoordinationAgent} variant receives a structured request and applies the same explicit predicates without LLM interpretation. The first baseline tests the contribution of symbolic authorization; the second tests the contribution of linguistic mediation. Mixed-agent scenarios must include remote A2A agents, local Python agents, and Lingo-backed linguistic agents so that all specializations are evaluated through the same SDK boundary.

Additional ablations remove one mechanism at a time: contract validation, explicit authority checks, trace recording, or bounded-auxiliary rules. These ablations are evaluated only where the removed mechanism is relevant to the scenario. They are not treated as independent general-purpose systems.

\section{Dimensiones y métricas de evaluación}

\begin{table}[htbp]
\centering
\caption{Dimensiones de evaluación, indicadores y evidencia.}
\label{tab:method-metrics}
\begin{tabular}{p{0.24\textwidth}p{0.30\textwidth}p{0.34\textwidth}}
\toprule
\textbf{Dimensión} & \textbf{Indicador} & \textbf{Métrica o evidencia} \\
\midrule
Juicio de viabilidad & Autorización o rechazo correcto & Exactitud, precisión, exhaustividad y F1 para la clasificación de planes viables \\
Validez del plan & Detección de candidatos inválidos & Tasa de rechazo de planes inválidos por categoría de fallo \\
Resultado de ejecución & Artefactos requeridos completados & Tasa de éxito de tareas y tasa de artefactos válidos por contrato \\
Trazabilidad & Evidencia requerida de decisión registrada & Completitud de campos de traza y tasa de decisiones reconstruibles \\
Eficiencia & Esfuerzo de coordinación & Latencia extremo a extremo, llamadas al planificador, mensajes y uso de tokens \\
Robustez & Estabilidad bajo variación controlada & Desempeño ante paráfrasis, pérdida de agentes y perturbación de metadatos \\
\bottomrule
\end{tabular}
\end{table}

Outcome metrics and process metrics are reported separately. A correct final artifact does not erase an unauthorized action or an invalid intermediate plan. Conversely, a justified refusal is a correct outcome for an infeasible request and must not be counted as task failure.

\section{Procedimiento experimental}

For each scenario, every applicable configuration receives semantically equivalent inputs and an identical admitted-agent snapshot. Stochastic configurations are executed repeatedly with recorded model, decoding parameters, prompt version, and random seed where supported. Each run stores the interpreted request, candidate plan, predicate outcomes, authorization decision, coordination-agent decisions, SDK method calls, agent kind, transport or runtime path, messages, status transitions, artifacts, validation results, latency, and resource use. Aggregate results are accompanied by per-category error analysis. The implementation version, coordination-agent version, SDK version, protocol version, hardware, and software environment must be fixed and reported before measurements are collected.

\section{Estrategia de análisis}

Feasibility decisions are compared with reference labels through a confusion matrix and category-level error rates. Paired scenario results are used for comparisons because configurations face the same requests. Descriptive statistics include central tendency and dispersion rather than averages alone. Where the sample size and distribution justify inference, paired statistical tests and effect sizes are reported. Qualitative analysis classifies false acceptance, false refusal, invalid decomposition, unsupported auxiliary synthesis, authority violation, artifact mismatch, and incomplete evidence.

\section{Validez y reproducibilidad}

Construct validity depends on whether the metrics capture justified coordination rather than surface task completion. This threat is mitigated by separating plan validity, execution success, communication-boundary compliance, coordination-ownership compliance, and trace completeness. Internal validity may be affected by unequal prompts, model nondeterminism, or baseline implementation quality; shared inputs, repeated runs, versioned configurations, and paired comparisons reduce these effects. External validity is limited by the selected scenarios, admitted agents, SDK adapters, and A2A-backed implementation choices. The theoretical framework is protocol-independent, but empirical findings may not generalize to decentralized authorization, physical robots, adversarial registries, or other protocols. Conclusion validity depends on scenario count and variance and must therefore be assessed after data collection.

Reproducibility requires publication of scenario specifications, reference labels, registry snapshots, configuration files, prompts, raw traces, analysis scripts, and exact dependency versions. Sensitive credentials and private agent endpoints are excluded, but their interfaces and response fixtures should be represented by reproducible substitutes.

\chapter{Proposed Theoretical Framework}
\label{chap:theoretical-solution}

\begin{figure}[htbp]
\centering
\begin{tikzpicture}[
  node distance=8mm and 8mm,
  every node/.style={font=\small},
  box/.style={draw, rounded corners, align=center, minimum height=8mm, inner sep=4pt},
  coord/.style={box, fill=blue!10, minimum width=42mm, minimum height=22mm},
  sym/.style={box, fill=blue!8},
  lng/.style={box, fill=orange!12},
  outc/.style={box, fill=green!10, minimum width=34mm},
  grp/.style={draw, dashed, rounded corners, inner sep=5pt},
  arr/.style={-Latex, thick},
  reg/.style={box, fill=gray!8, minimum width=28mm, minimum height=6mm}
]
\node[lng] (user) {User request\\(\textit{NL})};
\node[coord, right=18mm of user] (coord) {\textbf{Coordinator}\\[2mm]
  \begin{tabular}{cccc}
  \scriptsize Interpreter & \scriptsize Planner & \scriptsize Analyzer & \scriptsize Generator
  \end{tabular}};
\node[reg, above=14mm of coord] (reg) {Registered agents\\$K(R)$};
\node[outc, below left=10mm and 4mm of coord] (acc) {Accepted Plan $\Pi$};
\node[outc, below right=10mm and 4mm of coord] (inf) {Infeasibility $\bot$};

\draw[arr] (user) -- node[above]{\scriptsize T1} (coord);
\draw[arr] (reg) -- node[right]{\scriptsize capabilities} (coord.north west);
\draw[arr] (coord.south) -- ++(0,-2mm) -| (acc.north) node[pos=0.25,above]{\scriptsize feasible};
\draw[arr] (coord.south) -- ++(0,-2mm) -| (inf.north) node[pos=0.25,above]{\scriptsize $\lnot$ feasible};

\begin{scope}[on background layer]
  \node[grp, fit=(coord)(acc)(inf)(reg), label=above:\scriptsize Framework contribution] {};
  \node[grp, fit=(user), label=above:\scriptsize Linguistic surface] {};
\end{scope}
\end{tikzpicture}
\caption{End-to-end view of the proposed coordination process. A user request enters the system, is interpreted by the linguistic layer, decomposed and validated against the symbolic capability space, and is either executed as a coordinated plan of agent commitments or returned as an explicit infeasibility verdict. The coordinator is the central element; the user, the registry, and the two terminal outcomes are its environment.}
\label{fig:process-overview}
\end{figure}

\section{Purpose of the Framework}

The solution proposed in this thesis is a general model for coordinating heterogeneous agents in an open multi-agent system that bridges linguistic reasoning, provided by large language models, with symbolic coordination, expressed through explicit capabilities, plans, and commitments. Its purpose is not to prescribe a single software protocol, programming language, model provider, or deployment platform. Instead, it defines the conceptual conditions under which a group of autonomous agents can be treated as a coordinated problem-solving system. The framework assumes that agents may differ in internal architecture, reasoning method, language capability, and ownership, but that they can participate in a common coordination process if their capabilities and interaction boundaries are sufficiently described.

The central research problem is therefore one of controlled openness. A useful multi-agent system should be able to accept diverse agents, discover what they can do, determine whether a requested problem is solvable within the current system boundary, and coordinate a solution when feasible. At the same time, such a system should avoid unnecessary complexity. The decision process should remain as simple as possible: if a direct capability match is sufficient, no elaborate planning procedure is required; if a decomposition is necessary, it should be no deeper than the problem demands; if the system lacks a required capability, it should return a clear infeasibility explanation instead of simulating competence.

The thesis hypothesis is that such a system is achievable by combining a centrally coordinated capability-matching loop with a disciplined boundary between linguistic interpretation and symbolic execution. The framework formalizes this hypothesis as a small set of objects, a deterministic feasibility predicate, a bounded rule set for auxiliary capability synthesis, and a set of properties the loop is intended to satisfy.

\section{Assumptions and Scope}

The framework rests on a compact set of assumptions. First, participating agents are treated as black boxes with respect to their internal implementation. An agent may contain symbolic rules, a large language model, a classical planner, a retrieval system, a statistical model, a deterministic program, or a combination of these mechanisms. What is observable to the coordinator is the externally declared capability profile. Second, the coordinator is a single, trusted component that maintains the global view of the problem, the available capabilities, and the current solution plan. This is a deliberate departure from fully decentralized coordination in exchange for an explicit feasibility decision point. Third, capability descriptions are static within a coordination session and are provided honestly by participating agents. Fourth, the communication protocol layer satisfies the required protocol properties stated in Section~\ref{sec:protocol} (Protocol Substitutability). Fifth, the symbolic layer is assumed to be correct with respect to its own specifications: a deterministic feasibility check is trusted to return a sound verdict on the explicit constraints it examines.

The framework makes no claim about a number of related concerns and explicitly excludes them from scope. It does not verify the internals of language models, since the controllable boundary in this work is the symbolic envelope around the model rather than the model itself. It does not replace low-level agent communication languages such as KQML or FIPA-ACL \cite{finin1994kqml,labrou1997proposal,fipa2002acl}; such protocols may be used for transport and message structure. It does not provide a solution to decentralized consensus, Byzantine fault tolerance, or open-market task allocation in adversarial settings. It does not guarantee optimality of the resulting coordination plan, only that the plan is justified relative to the declared capabilities. Finally, it does not claim that generated auxiliary specifications are equivalent in quality to dedicated, long-trained agents; they are narrow, validated substitutes used when justified by simplicity and boundedness.

\section{Design Requirements for the MAS}

The theoretical framework is guided by a compact set of design requirements. First, the system must be open: agents should be able to join, leave, or become unavailable without requiring a redesign of the entire architecture. Openness does not imply unrestricted participation; it requires explicit admission, description, and trust boundaries. Second, the system must be vendor independent. The coordination model should not depend on a particular LLM provider, orchestration framework, cloud platform, or proprietary runtime. A participating agent should be evaluated by its declared capabilities and observable behavior rather than by the technology used internally.

Third, the system must be interoperable. Agents require shared conventions for describing capabilities, exchanging messages, returning artifacts, and reporting errors. Fourth, the system is distributed by design. Agents remain independently deployable units, and coordination occurs through explicit interactions rather than shared internal state. Fifth, the system must preserve coordination. Although agents are distributed, the solution process requires a central coordinating role that maintains the global view of the problem, the available capabilities, and the current solution plan.

Sixth, the system must prefer simplicity whenever possible. The coordinator should use the least complex decision rule that can justify a correct action. A single-agent direct solution is preferable to a multi-agent workflow when the former is sufficient. A deterministic capability check is preferable to an LLM judgment when the condition is explicit. A generated auxiliary agent should be proposed only when it reduces complexity or fills a narrow missing language capability. Seventh, the system must support a symbolic execution envelope: plans, commitments, and tool calls must be expressible in machine-checkable form so that they can be validated before execution. Finally, the system must be observable and bounded: decisions, delegations, failures, and produced artifacts should be traceable, and agents should operate under explicit authority boundaries.

\section{Agents and Capabilities}

In the proposed framework, an agent is an autonomous computational participant that can accept tasks, perform actions, and return results within a declared boundary of competence. The internal implementation of the agent is intentionally abstract. What matters for coordination is not the internal mechanism alone, but the externally visible capability profile.

\textbf{Definition 3.1 (Agent).} An agent is a tuple $a = (\mathit{id}, C(a), M_{\mathrm{in}}(a), M_{\mathrm{out}}(a), \tau(a))$ where $\mathit{id}$ is a unique identifier, $C(a)$ is a finite set of declared capabilities, $M_{\mathrm{in}}(a)$ and $M_{\mathrm{out}}(a)$ are finite sets of accepted input and produced output modalities, and $\tau(a)$ is a trust level drawn from a partial order.

A capability is a declared ability to transform certain inputs into certain outputs under known constraints. Capabilities may be functional, such as generating a plan, classifying a document, querying a database, validating a formal expression, or producing an artifact. They may also be communicative, such as explaining a result, negotiating a dependency, or translating between terminology used by different agents. For a coordinator, a capability description is useful only if it supports matching between problem requirements and agent responsibilities.

\textbf{Definition 3.2 (Capability).} A capability is a tuple $c = (\mathit{name}, \Sigma_{\mathrm{in}}, \Sigma_{\mathrm{out}}, \gamma, \nu)$ where $\Sigma_{\mathrm{in}}$ and $\Sigma_{\mathrm{out}}$ are input and output schemas, $\gamma$ is a set of explicit constraints (preconditions, authority limits, resource bounds), and $\nu$ is a deterministic validator that decides whether a candidate output satisfies the capability's contract.

The capability space is the abstract set of all capabilities currently available to the system. The system boundary is defined by this space: a task is feasible only if its required capabilities can be mapped to existing agents or to acceptable auxiliary capabilities synthesized for the solution. This formulation prevents the coordinator from treating every problem as solvable. The system may be intelligent and adaptive, but it remains bounded by what its participating agents can actually perform.

\textbf{Definition 3.3 (Capability Space and System Boundary).} Given a registry $R$ of currently available agents, the capability space is $K(R) = \bigcup_{a \in R} C(a)$. The system boundary is the pair $\beta(R) = (K(R), M(R))$ where $M(R) = \bigcup_{a \in R} M_{\mathrm{in}}(a) \cup M_{\mathrm{out}}(a)$ is the set of supported modalities.

\section{Linguistic Agents}

A linguistic agent is an agent whose primary contribution is language-mediated processing. Its tasks may include interpretation, summarization, classification, translation, information extraction, semantic matching, question answering, report generation, or natural-language explanation. Linguistic agents are necessary because many coordination problems enter the system in natural language, while many system actions require structured interpretation. The linguistic layer therefore mediates between human intent, agent capability descriptions, symbolic constraints, and final explanations.

A linguistic agent is not necessarily a full LLM agent. Some linguistic tasks require broad reasoning, context-sensitive interpretation, or planning over ambiguous instructions; these may justify an LLM-based agent. Other tasks are simpler and may be handled by a smaller NLP model, a rule-based extractor, a retrieval component, a classifier, or a deterministic template. For example, language detection, keyword extraction, entity recognition, intent classification, or conversion of a narrow text pattern into a structured field may not require a general-purpose LLM. Treating all language work as LLM work would increase cost, latency, and complexity without necessarily improving reliability.

The framework therefore maps linguistic agents to capabilities rather than to one implementation class. A linguistic capability should describe the kind of language operation performed, the accepted input forms, the expected output artifacts, the limits of confidence, and the conditions under which the result should be validated by another agent. This capability-oriented view keeps the theory general. It allows the system to include powerful LLM-based agents, lightweight NLP agents, and hybrid language-processing agents under the same conceptual model.

\section{The Linguistic--Symbolic Bridge}

A central component of the framework is the explicit boundary between linguistic reasoning and symbolic execution. The bridge is not a single translator but a set of four named transitions, each with a producer, a consumer, a data shape, a trust assumption, and a failure mode.

\textbf{Transition 1 (NL $\rightarrow$ Symbolic Problem).} The user problem, expressed in natural language, is transformed by an LLM-based interpreter into a structured $\mathit{ProblemRequest}$ (see Definition 3.4). The interpreter is trusted to produce a faithful structured representation, but the symbolic layer does not depend on the natural-language text thereafter. Failure mode: the interpreter misrepresents the user intent; mitigation is to expose the resulting $\mathit{ProblemRequest}$ to the user or a verification step.

\textbf{Transition 2 (Symbolic Problem $\rightarrow$ Candidate Plan).} Given the $\mathit{ProblemRequest}$ and the current capability index, an LLM-based planner proposes a candidate coordination plan, expressed in symbolic form. The planner is trusted to propose, not to decide. The plan is checked deterministically before it is accepted. Failure mode: the planner invents a capability that does not exist; mitigation is the feasibility analyzer (see Section~\ref{sec:feasibility-boundary}).

\textbf{Transition 3 (Symbolic Plan $\rightarrow$ NL Explanation).} A symbolic plan is rendered back into natural language for human review, for follow-up questions, or for delegation contexts where another agent expects a prose brief. The renderer is trusted only for presentation; the underlying plan is the authoritative artifact. Failure mode: a misleading explanation; mitigation is to keep the plan and the explanation as separate objects linked by identifier.

\textbf{Transition 4 (Observation $\rightarrow$ Structured Feedback).} External observations, including tool errors, agent failures, infeasibility reports, and protocol violations, are transformed into structured feedback that updates the plan state. This transition may be performed by a small classifier, a rule, or an LLM, depending on the observation's structure. The symbolic layer is trusted to update only the parts of the plan that the feedback explicitly references.

The bridge aligns with classical agent theory. The BDI separation between beliefs, desires, and intentions \cite{rao1991modeling,rao1995bdi} corresponds to the distinction between the structured problem (belief), the candidate plan (desire), and the accepted plan (intention). ACL performatives such as request, inform, and commit \cite{austin1962how,searle1969speech,cohen1990intention,fipa2002acl} correspond to message-level operations on the plan. The framework uses the bridge as the architectural embodiment of this discipline: the LLM produces, the symbolic layer decides.

\begin{figure}[htbp]
\centering
\begin{tikzpicture}[
  node distance=10mm and 6mm,
  every node/.style={font=\small},
  box/.style={draw, rounded corners, align=center, minimum height=8mm, minimum width=24mm, inner sep=3pt},
  sym/.style={box, fill=blue!8},
  lng/.style={box, fill=orange!12},
  grp/.style={draw, dashed, rounded corners, inner sep=4pt},
  arr/.style={-Latex, thick}
]
\node[lng] (user) {User problem\\(\textit{NL})};
\node[sym, right=of user] (req) {ProblemRequest\\$P = (g, R, \mathit{ctx})$};
\node[lng, right=of req] (plan) {Candidate plan\\via LLM Planner};
\node[sym, right=of plan] (feas) {Feasibility\\Analyzer};
\node[sym, below=of feas] (aux) {Generated $\mathit{aux}$\\if synthesizable};
\node[sym, left=of aux] (capidx) {Capability\\index $K(R)$};
\node[sym, right=of aux] (acc) {Accepted plan\\$\Pi$ or $\bot$};
\node[sym, below=of acc] (comm) {Commitments\\$\mathcal{C}(\Pi)$};
\node[sym, left=of comm] (proto) {Protocol\\transport};
\node[sym, left=of proto] (agg) {Artifact\\Aggregator};
\node[lng, left=of agg] (out) {Result or\\infeasibility};

\draw[arr] (user) -- node[above]{\scriptsize T1} (req);
\draw[arr] (req) -- node[above]{\scriptsize T2} (plan);
\draw[arr] (plan) -- (feas);
\draw[arr] (capidx) -- (feas);
\draw[arr] (feas) -- node[above,sloped]{\scriptsize aux if gap} (aux);
\draw[arr] (aux) -- (acc);
\draw[arr] (feas) -- node[above]{\scriptsize feasible} (acc);
\draw[arr] (acc) -- (comm);
\draw[arr] (comm) -- (proto);
\draw[arr] (proto) -- (agg);
\draw[arr] (agg) -- (out);

\begin{scope}[on background layer]
  \node[grp, fit=(req)(acc)(capidx)(aux), label=above:\scriptsize Symbolic layer] {};
  \node[grp, fit=(user)(plan), label=above:\scriptsize Linguistic layer] {};
  \node[grp, fit=(comm)(proto)(agg)(out), label=below:\scriptsize Execution and feedback] {};
\end{scope}
\end{tikzpicture}
\caption{Data flow across the linguistic--symbolic bridge. Transitions T1 (NL $\rightarrow$ ProblemRequest) and T2 (ProblemRequest $\rightarrow$ candidate plan) are produced by the LLM planner; the remaining transitions are deterministic on the symbolic side. The bridge ensures that no commitment is created without a corresponding plan and that no plan is accepted without feasibility verification.}
\label{fig:theoretical-bridge}
\end{figure}

\section{Auxiliary Linguistic Capability Synthesis}

The coordinator may encounter a problem whose decomposition requires a simple linguistic capability that is not currently represented by an existing agent. In such cases, the framework allows auxiliary linguistic capability synthesis. This means that the coordinator may propose a temporary, task-local linguistic agent as part of the solution plan. The generated agent is not assumed to become a permanent member of the system. It is a bounded capability specification created to satisfy a narrow subtask.

This concept is important because it preserves both openness and simplicity. If a problem requires only a small text-normalization step, it may be excessive to delegate to a full LLM agent or to reject the whole problem because no registered agent advertises that exact micro-capability. A synthesized auxiliary linguistic agent can fill the gap if its purpose, inputs, outputs, validation rule, and lifecycle are explicit. The coordinator should not silently invent complex agents. It should expose the generated capability as part of the solution proposal so that the system can evaluate whether the addition is acceptable.

\textbf{Definition 3.4 (Generated Auxiliary Specification).} An auxiliary specification is a tuple $\mathit{aux} = (\mathit{purpose}, \Sigma_{\mathrm{in}}, \Sigma_{\mathrm{out}}, \mathit{method}, \nu, \lambda)$ where $\mathit{purpose}$ is a short human-readable description, $\Sigma_{\mathrm{in}}$ and $\Sigma_{\mathrm{out}}$ are schemas, $\mathit{method}$ is an enumerated or closed-form procedure, $\nu$ is a deterministic validator, and $\lambda$ is a lifecycle token that bounds the validity of $\mathit{aux}$ to the current plan.

Auxiliary synthesis is constrained by four explicit rules.

\textbf{Rule A1 (Closed method).} Synthesis is allowed only for capabilities whose $\mathit{method}$ can be enumerated, expressed as a closed-form procedure, or specified as a small set of deterministic steps. Open-ended language generation does not qualify.

\textbf{Rule A2 (Validator present).} A valid $\nu$ must be specified at synthesis time. Capabilities without a deterministic validator cannot be synthesized.

\textbf{Rule A3 (No side effects).} Synthesis is forbidden when the gap requires persistent state writes, external authority, or privileged access beyond the coordinator's validation capacity.

\textbf{Rule A4 (Lifecycle bound).} Every synthesized $\mathit{aux}$ is created with a lifecycle token $\lambda$ that ties its validity to the current coordination plan. It does not enter the registry automatically.

For example, a decomposition that requires extracting email addresses from a short free-text field is an acceptable synthesis target: the method is a regular expression, the validator confirms the output matches the email schema, and the lifecycle is bound to the current plan. A decomposition that requires medical-image classification is not an acceptable target: the method is not enumerable, the validator cannot be specified in advance, and the capability is not bounded by the current plan.

\section{Inter-Agent Communication Protocol}
\label{sec:protocol}

The framework treats the communication protocol as an infrastructure layer rather than as a component of its symbolic reasoning. The protocol carries task delegations, capability descriptions, status updates, and artifacts between the coordinator and participating agents. Its responsibilities are limited to transport, addressing, and delivery; it does not interpret the content of the messages, decide feasibility, or expose agent internals. Because the framework does not depend on a particular protocol, the implementation chapter is free to select or compose any protocol that satisfies the required properties stated below. The framework remains sound as long as the protocol in use satisfies these properties, regardless of its name or its vendor.

The following eight properties describe what the framework requires of the protocol. They are stated as design intentions; the implementation chapter is responsible for demonstrating that the chosen protocol satisfies them or for showing how a composition of protocols can satisfy the full set.

\textbf{Property 3.5 (Addressability).} Every participating agent has a stable, unique identifier reachable through the protocol. The coordinator can address any registered agent by this identifier and can confirm its liveness before delegating a task.

\textbf{Property 3.6 (Capability Advertisement).} The protocol can carry a structured capability description corresponding to Definition~3.2: input and output schemas, explicit constraints, and a deterministic validator. The coordinator uses these descriptions to match problem requirements against the available capability space $K(R)$ (Definition~3.3).

\textbf{Property 3.7 (Task Delegation).} The protocol can carry a typed delegation from the coordinator to a specific agent, comprising a task body, an expected artifact description, and a lifecycle token $\lambda$ (the same token used in Definition~3.4 for auxiliary specifications). The delegation is addressed to a single agent and is not broadcast.

\textbf{Property 3.8 (Status and Result Delivery).} The protocol can carry status updates (pending, running, completed, failed) and the resulting artifact from a delegated agent back to the coordinator. Status transitions are delivered in the order they occur, and the resulting artifact, if any, is delivered with an associated completion signal.

\textbf{Property 3.9 (Authenticated Origin).} Each message identifies its sender, and the protocol layer rejects messages whose origin cannot be verified against the trust assignment $\tau(a)$ (Definition~3.1) of the claimed agent. The coordinator treats an unauthenticated message as if the corresponding agent were unavailable.

\textbf{Property 3.10 (Ordered Delivery within a Session).} Messages belonging to the same coordination session are delivered in the order they were produced. Out-of-order delivery is detected and reported to the coordinator as a protocol violation, which the symbolic layer records as a structured feedback event under Transition~4 of the linguistic--symbolic bridge.

\textbf{Property 3.11 (Bounded Latency and Liveness).} Within a session, a delegated task receives a status update or a result within a declared time bound. If the bound is exceeded, the coordinator observes the absence of a reply and triggers an infeasibility response rather than blocking indefinitely. Liveness is therefore a property of the protocol as observed by the coordinator, not an internal guarantee of the participating agents.

\textbf{Property 3.12 (Independence from Internal Agent Architecture).} The protocol does not require agents to expose their internal reasoning, language model, memory, or tool registry. Only the artifacts listed in Properties 3.5--3.8 are visible across the protocol boundary. This property is what makes the framework vendor independent at the protocol level, in the same way that Definition~3.1 makes the framework vendor independent at the agent level.

\textbf{Property 3.13 (Protocol Substitutability).} Any protocol satisfying Properties 3.5--3.12 can be substituted for the framework's communication layer without affecting Properties 3.1--3.4 (soundness, simplicity, bounded auxiliary synthesis, observability). The framework's correctness therefore does not depend on a specific protocol choice. The implementation chapter is responsible for selecting or composing a protocol that satisfies the full set; whether the chosen protocol is widely deployed, open, or proprietary is an implementation concern rather than a theoretical one.

\section{Hierarchical Coordination and Commitments}

The framework uses hierarchical coordination in a limited and explicit sense. The central coordinator has authority over task interpretation, capability matching, decomposition, feasibility assessment, and orchestration of the solution plan. It does not erase the autonomy of participating agents. Each agent remains responsible for executing tasks within its declared boundary and for returning results or errors. The coordinator manages the global problem-solving process; agents perform bounded contributions.

This hierarchy is necessary because open distributed systems need a global decision point for feasibility. Without such a role, no component has responsibility for deciding whether the system as a whole can solve a problem. A user may request a result that requires several agents, but each individual agent may see only its local task. The coordinator maintains the overall context, determines dependencies, selects agents, and verifies that the combination of agents can plausibly satisfy the request.

The framework also introduces an explicit notion of commitment to address a gap repeatedly noted in the background: dialogue alone does not guarantee accountability \cite{singh1999commitments}. A commitment is created when the coordinator delegates a subtask to a specific agent. It records the debtor (the agent), the creditor (the coordinator on behalf of the user), the action to be performed, the expected artifact, and the conditions under which the commitment is discharged. The commitment persists in symbolic form regardless of what the participating agents say to each other in natural language. This makes it possible to detect, after the fact, whether an agent accepted a task, whether it executed, and whether the result satisfies the postcondition.

The coordinator's reasoning has two complementary modes. Semantic reasoning interprets the user problem and compares it with natural-language capability descriptions. Deterministic reasoning checks explicit constraints such as required input and output forms, permission boundaries, unavailable agents, missing capabilities, or incompatible dependencies. The theoretical framework does not require a particular algorithm for semantic reasoning. It requires only that semantic interpretation be bounded by explicit checks before a solution is accepted as feasible.

\section{Feasibility Boundary}
\label{sec:feasibility-boundary}

\textbf{Definition 3.5 (Problem Request).} A problem request is a tuple $P = (g, R, \mathit{ctx})$ where $g$ is a natural-language goal, $R$ is a set of required capability descriptors, and $\mathit{ctx}$ is auxiliary context. $R$ may include auxiliary-eligible descriptors, which the coordinator may satisfy by synthesis.

\textbf{Definition 3.6 (Coordination Plan).} A coordination plan is a tuple $\Pi = (T, \alpha, \prec, A_{\mathrm{out}}, \kappa)$ where $T$ is a finite set of subtasks, $\alpha : T \rightarrow \mathcal{A} \cup \mathcal{X}$ assigns each subtask to a registered agent or an auxiliary specification, $\prec$ is a partial order on $T$ expressing dependencies, $A_{\mathrm{out}}$ is a set of expected artifacts, and $\kappa$ is the completion criterion.

\textbf{Definition 3.7 (Feasibility).} Given a problem $P = (g, R, \mathit{ctx})$ and a registry $R$ (with capability space $K(R)$ and boundary $\beta(R)$), the feasibility predicate is
\[
\mathit{feasible}(P, R) \;\Longleftrightarrow\;
\begin{array}{l}
\forall c \in R : c \in K(R) \lor \mathit{synthesizable}(c), \\
\forall c \in R : \Sigma_{\mathrm{out}}(c) \subseteq M(R) \lor \text{bridgeable}, \\
\exists \Pi : T = R \land \alpha \text{ is total} \land \prec \text{ is a valid partial order}, \\
\forall t \in T : \tau(\alpha(t)) \geq \tau_{\mathrm{req}}(t).
\end{array}
\]

In words, a problem is feasible if every required capability is either available in the registry or synthesizable under rules A1--A4, if all required input and output modalities are present or can be bridged, if a valid coordination plan can be constructed with a total assignment and a consistent partial order, and if every assigned agent meets the trust requirement of its subtask. A problem is infeasible if any of these conditions fail. The infeasibility response identifies the specific failure: a missing capability, an unbridgeable modality, an unsatisfiable dependency, or an authority gap.

The feasibility boundary is not fixed permanently. It changes as agents join, leave, update their capabilities, or become unavailable. It also changes when the coordinator is allowed to synthesize simple auxiliary linguistic capabilities. However, the boundary must always be explicit at the moment a problem is evaluated. The system should avoid pretending to solve tasks outside its boundary. A principled error is preferable to an unsupported plan.

\section{Properties of the Framework}

The framework is intended to satisfy a small set of named properties. No formal proofs are given; the statements are intended to be testable in the implementation and results chapters.

\textbf{Property 3.1 (Soundness of Feasibility).} If the coordinator returns $\Pi$ as feasible, then for every subtask $t \in T$, the assigned executor $\alpha(t)$ is either a registered agent whose declared capability covers the subtask's requirements or a synthesized auxiliary specification that passes rules A1--A4. The system does not return a plan it cannot justify.

\textbf{Property 3.2 (Simplicity Preference).} If a problem $P$ can be satisfied by direct delegation to a single agent $a$, the coordinator returns a plan of size one. Decomposition is introduced only when no single agent covers $R$.

\textbf{Property 3.3 (Bounded Auxiliary Synthesis).} The set of auxiliary specifications produced for a plan $\Pi$ is finite, each member satisfies A1--A4, and no auxiliary specification persists beyond the lifecycle token $\lambda$ attached to its plan.

\textbf{Property 3.4 (Observability).} Every accepted plan $\Pi$ is stored with the trace of its decisions: which agent or auxiliary was selected for each subtask, which constraints were checked, which infeasibility reasons were considered, and which execution events occurred. The trace is sufficient to distinguish agent failure from coordination failure after the fact.

These properties are stated as design intentions rather than theorems. They provide the criteria against which the implementation and the empirical evaluation in Chapter~\ref{chap:results} can be assessed.

\section{Worked Example}

Consider a small registry of three agents, each advertising a single capability.

\begin{itemize}
  \item $a_1$ (Summarizer): summarization of English text, input modality text, output modality text.
  \item $a_2$ (Calculator): arithmetic evaluation of an expression in a fixed grammar, input modality $\mathit{expr}$, output modality $\mathit{number}$.
  \item $a_3$ (DB-Query): execution of a read-only SQL query against a registered database, input modality $\mathit{sql}$, output modality $\mathit{table}$.
\end{itemize}

A user submits the following natural-language request: ``Summarize the quarterly report, then compute the percentage change in revenue between Q1 and Q2, and finally list the top five customers by total order value.''

The LLM interpreter produces the structured $\mathit{ProblemRequest}$ $P = (g, R, \mathit{ctx})$ with $R = \{\mathit{summarize}, \mathit{arithmetic}, \mathit{top\_k\_sql}\}$. The capability matcher finds direct coverage for $\mathit{summarize}$ (via $a_1$) and $\mathit{arithmetic}$ (via $a_2$). The $\mathit{top\_k\_sql}$ capability is not present; however, the coordinator recognizes that $a_3$ accepts a SQL query and that the missing piece is a small translation step from "top five customers by total order value" into a SQL string. The coordinator generates an auxiliary specification $\mathit{aux} = (\mathit{sql\_templating}, \Sigma_{\mathrm{in}} = \text{schema}, \Sigma_{\mathrm{out}} = \mathit{sql}, \mathit{method} = \text{template}, \nu = \text{parse-and-validate}, \lambda = \Pi_1)$. The method is enumerable, the validator parses the generated SQL against the database schema, and the lifecycle is bound to the current plan.

The candidate plan is $\Pi_1 = (\{t_1, t_2, t_3\}, \alpha, \prec, A_{\mathrm{out}}, \kappa)$ with $\alpha(t_1) = a_1$, $\alpha(t_2) = a_2$, $\alpha(t_3) = \mathrm{exec}(a_3, \mathit{aux})$, $\prec = \{t_1 \prec t_3, t_2 \prec t_3\}$, $A_{\mathrm{out}} = \{\mathit{summary}, \mathit{pct}, \mathit{table}\}$, and $\kappa$ requiring all three artifacts to be present and well-formed. The Feasibility Analyzer confirms coverage, modality compatibility, total assignment, partial-order consistency, and trust levels. The plan is approved.

If, instead, the user had requested "summarize the quarterly report and predict next year's revenue," the request would be infeasible: the second subtask requires a forecasting capability that is neither registered nor synthesizable under rules A1--A4. The coordinator returns a structured infeasibility response identifying the missing capability and the reason it cannot be synthesized.

Figure~\ref{fig:process-with-example} instantiates the process overview for the present example. The three agents supply the capability index, the generator produces $\mathit{aux}$ under rules A1--A4, and the outcome follows the path determined by the feasibility predicate.

\begin{figure}[htbp]
\centering
\begin{tikzpicture}[
  node distance=8mm and 8mm,
  every node/.style={font=\small},
  box/.style={draw, rounded corners, align=center, minimum height=8mm, inner sep=4pt},
  coord/.style={box, fill=blue!10, minimum width=46mm, minimum height=28mm},
  sym/.style={box, fill=blue!8},
  lng/.style={box, fill=orange!12},
  outc/.style={box, fill=green!10, minimum width=36mm},
  arr/.style={-Latex, thick},
  reg/.style={box, fill=gray!8, minimum width=24mm, minimum height=6mm}
]
\node[lng] (user) {User request\\\scriptsize ``Summarize, compute pct, top 5.''};
\node[coord, right=18mm of user] (coord) {\textbf{Coordinator}\\[2mm]
  \begin{tabular}{c}
  \scriptsize T1: $\mathrm{NL} \to \mathit{ProblemRequest}$ \\
  \scriptsize T2: candidate plan $\Pi_1$ \\
  \scriptsize Feasibility check (Def.~3.7) \\
  \scriptsize Generator: $\mathit{aux}$ if gap
  \end{tabular}};
\node[reg, above=14mm of coord] (reg) {Registered agents};
\node[reg, below=2mm of reg] (a1) {$a_1$ Summarizer};
\node[reg, below=1mm of a1] (a2) {$a_2$ Calculator};
\node[reg, below=1mm of a2] (a3) {$a_3$ DB-Query};
\node[outc, below left=10mm and 4mm of coord] (acc) {Accepted $\Pi_1$};
\node[sym, below=2mm of acc, minimum width=36mm] (comm) {Commitments $\to$\\Protocol transport};
\node[outc, below right=10mm and 4mm of coord] (inf) {Infeasibility $\bot$};

\draw[arr] (user) -- node[above]{\scriptsize T1} (coord);
\draw[arr] (reg) -- node[right]{\scriptsize $K(R)$} (coord.north west);
\draw[arr] (coord.south) -- ++(0,-2mm) -| (acc.north) node[pos=0.25,above]{\scriptsize feasible};
\draw[arr] (coord.south) -- ++(0,-2mm) -| (inf.north) node[pos=0.25,above]{\scriptsize $\lnot$ feasible};
\draw[arr] (acc) -- (comm);
\end{tikzpicture}
\caption{Worked example instantiated against the process overview. The three agents supply the capability index, the generator produces the auxiliary spec $\mathit{aux}$ under rules A1--A4, and the outcome follows the path determined by the feasibility predicate.}
\label{fig:process-with-example}
\end{figure}

\section{Scope, Limitations, and Open Questions}

The framework is intentionally narrow. It assumes honest capability descriptions, static capabilities within a session, and a single trusted coordinator. It does not address the design of the underlying LLM, the selection of training data, or the optimization of model parameters. It does not provide a fully decentralized coordination algorithm, and it does not provide economic mechanisms for resource allocation in competitive settings. It also assumes that the deterministic checks it performs are themselves correct with respect to the symbolic specifications they examine; a faulty feasibility analyzer would compromise Property~3.1.

Several questions remain open. First, capability descriptions are currently static; future work could allow agents to refine their declared capabilities based on observed performance, while preserving the honesty assumption. Second, the bridge between linguistic and symbolic representations could be made more robust by requiring a verification step on the structured $\mathit{ProblemRequest}$ before downstream planning. Third, commitments are recorded as symbolic artifacts; richer commitment protocols, including conditional commitments and nested commitments, could extend the framework to more complex multi-agent negotiations. Fourth, the framework currently relies on a single coordinator; exploring how the same feasibility discipline can be preserved under decentralized or hierarchical-multi-coordinator settings is a natural next step.

The framework therefore defines a small, defensible contribution: a set of objects and a feasibility discipline that separates linguistic interpretation from symbolic execution, treats infeasibility as a normal outcome, and bounds the conditions under which auxiliary capabilities may be synthesized. The next chapter instantiates this framework through a concrete agent communication protocol and software design.

\chapter{Diseño del sistema: agente de coordinación sobre middleware SDK}
\label{chap:implementation-design}

\section{Instanciación del marco}

This chapter instantiates the theoretical framework from Chapter~\ref{chap:theoretical-solution} as a Python coordination agent built over a middleware SDK. The implementation design keeps the theoretical contribution protocol-independent, but it chooses the Agent2Agent (A2A) protocol as the canonical interoperability backbone for concrete communication. A2A is an open protocol intended to support communication and coordination among agents built with different technologies. Its project documentation describes standardized communication, Agent Cards for discovery, flexible interaction patterns, and rich data exchange through text, files, and structured JSON \cite{a2aproject2025}. Google's announcement similarly frames A2A as an open protocol for secure information exchange and coordination among agents across enterprise platforms \cite{google2025a2a}.

The implementation separates two roles. \texttt{CoordinationSdk} is the communication and runtime middleware. It registers remote A2A agents, local Python agents, and Lingo-backed linguistic agents; normalizes their capabilities; dispatches authorized tasks; converts artifacts; records runtime failures; and exposes communication traces. It does not plan, decide feasibility, authorize decompositions, or coordinate complete user requests.

\texttt{CoordinationAgent} is the central coordinating role. It receives the user request, uses the SDK to inspect the registry and capability index, invokes SDK-registered linguistic capabilities when interpretation or candidate planning is needed, builds direct or decomposed candidate plans, calls the symbolic \texttt{FeasibilityAnalyzer}, authorizes or refuses the plan, dispatches only authorized tasks through the SDK, and aggregates the final result. The agent never calls raw A2A transports, raw local handlers, or raw Lingo runtimes directly. All execution traffic still passes through \texttt{CoordinationSdk}.

A2A remains the default wire model for independently deployed agents. Remote agents are admitted through A2A Agent Cards and invoked through A2A-compatible task exchange. Local Python agents and Lingo-based linguistic agents are registered through the SDK and normalized into the same internal agent records, with identity, capabilities, input and output modes, trust metadata, invocation endpoint, and validation contract. Local and linguistic agent authors are therefore not required to implement A2A plumbing, while the coordination agent can reason over a uniform registry.

The prototype also uses Lingo as an optional linguistic and context-engineering substrate \cite{lingo2026}. Lingo is not hidden inside the SDK as a coordinator. Instead, Lingo-backed interpretation, extraction, classification, or candidate-planning behavior is exposed as SDK-registered linguistic capabilities. Non-linguistic local agents can use the SDK without depending on Lingo. The prototype dependency set pins \texttt{lingo-ai==2.0.2} and uses \texttt{a2a-sdk}, \texttt{pydantic}, \texttt{httpx}, \texttt{pytest}, and \texttt{pytest-asyncio} for the executable implementation and tests.

Lingo is therefore an implementation witness for the abstract transitions defined in Chapter~\ref{chap:theoretical-solution}. A Lingo \texttt{Context} can serve as the mutable linguistic ledger used by interpretation and trace generation. \texttt{Engine.create}, \texttt{Engine.decide}, and \texttt{Engine.choose} can implement controlled proposal operations. \texttt{Flow} can encode procedural candidate-planning workflows, while \texttt{StateMachine} can represent deterministic workflow states. Pydantic-backed schemas can materialize capability and artifact contracts, and \texttt{@app.when} reflexes can provide early guardrails or context injection. These mechanisms support interpretation, proposal, and feedback grounding; they do not authorize execution by themselves.

The mapping between formal objects and implementation records is direct. The records defined in Section~\ref{sec:data-structures} materialize the formal objects of Chapter~\ref{chap:theoretical-solution} as follows.

\begin{itemize}
  \item \texttt{CoordinationAgent} materializes the central coordinating authority. It owns interpretation workflow, direct matching, decomposition, feasibility checks, auxiliary decisions, authorization, and final result construction.
  \item \texttt{CoordinationSdk} materializes the communication abstraction. It owns registry admission, A2A gateway behavior, local and Lingo runtimes, task dispatch, artifact normalization, runtime failure recording, and communication trace access.
  \item \texttt{AgentRegistryEntry} materializes an agent $a$ (Definition~\ref{def:agent}), supplying $\mathit{id}$, declared capabilities $C(a)$, input modes $M_{\mathrm{in}}(a)$, output modes $M_{\mathrm{out}}(a)$, local admission policy $\tau(a)$, observable status $\sigma(a)$, and the invocation endpoint used by the SDK.
  \item \texttt{CapabilityRequirement} materializes a required capability descriptor $q \in \mathcal{Q}$ with modalities, schemas, constraints, trust level, and auxiliary eligibility.
  \item \texttt{ProblemRequest} materializes the request $P = (g, \mathcal{Q}, \chi, \mathcal{O}, \Gamma)$ from Definition~\ref{def:problem-request}.
  \item \texttt{SolutionProposal} materializes the coordination plan $\Pi$ from Definition~\ref{def:coordination-plan}, including task requirements, assignments, dependencies, auxiliary specifications, expected artifacts, and completion criteria.
  \item \texttt{FeasibilityReport} records the verdict and predicate evidence from Definition~\ref{def:feasibility}. It is produced by the \texttt{FeasibilityAnalyzer} at the \texttt{CoordinationAgent} layer.
  \item \texttt{GeneratedNlpAgentSpec} materializes the auxiliary capability specification from Definition~\ref{def:auxiliary}, bounded by Rules~\ref{rule:approved-construction}--\ref{rule:lifecycle}.
  \item \texttt{CoordinatorState} is the coordination agent's session state containing the registry snapshot, interpreted request, candidate plan, feasibility evidence, and trace.
  \item \texttt{TraceEvent} records auditable observations across interpretation, authorization, delegation, and artifact handling.
\end{itemize}

Each of the five properties stated in Chapter~\ref{chap:theoretical-solution} (justified acceptance, simplicity, bounded auxiliary synthesis, traceability, and explicit refusal) corresponds to a concrete decision or trace event owned by the \texttt{CoordinationAgent}, supported by communication evidence supplied through the SDK.

\section{Componentes del sistema}

The proposed implementation is organized around a coordinating agent above a middleware SDK. The \texttt{CoordinationAgent} is the planner and authorizer. It receives user requests, maintains coordination state, requests linguistic interpretation when needed, constructs candidate plans, applies deterministic feasibility checks, decides whether to authorize or refuse the plan, dispatches authorized work through the SDK, and aggregates the final result.

\texttt{CoordinationSdk} is the communication and runtime substrate. It is the facade through which agents are registered and invoked. The SDK may use private helper classes internally, but adapter inheritance and transport details are not part of the public API. The SDK can reject invalid configuration, malformed registrations, missing handlers, invalid Agent Cards, and transport-level failures, but it does not decide whether a user request is feasible as a multi-agent plan.

The private SDK subsystems keep communication responsibilities simple. The Registry and Capability Index maintain admitted agent records and searchable capability metadata. The A2A Gateway fetches and validates Agent Cards, sends A2A tasks, and normalizes remote status, artifacts, cancellations, timeouts, and failures. The Local Agent Runtime wraps Python functions, scripts, classes, symbolic planners, and deterministic tools. The Lingo Agent Runtime wraps Lingo apps, flows, tools, state machines, and guardrails when linguistic behavior is required. The SDK Trace Store records communication and runtime observations. The Artifact Normalizer converts outputs into validator-friendly records.

\begin{table}[htbp]
\centering
\caption{Componentes de implementación mapeados a conceptos teóricos.}
\label{tab:implementation-components}
\begin{tabular}{p{0.30\textwidth}p{0.58\textwidth}}
\toprule
\textbf{Componente} & \textbf{Rol en la implementación} \\
\midrule
\texttt{CoordinationAgent} & Posee el flujo de interpretación, correspondencia directa, descomposición, verificaciones de viabilidad, decisiones auxiliares, autorización y construcción del resultado final. \\
\texttt{CoordinationSdk} & Expone operaciones de middleware para registro, inspección del registro, indexación de capacidades, despacho de tareas, normalización de artefactos, reporte de fallos en tiempo de ejecución, trazas y reinicio de sesión. \\
\texttt{FeasibilityAnalyzer} & Aplica verificaciones deterministas para \texttt{CoordinationAgent}; no forma parte de la frontera de comunicación del SDK. \\
Pasarela A2A & Gestiona la ingesta de Agent Cards y la comunicación compatible con A2A con agentes remotos. \\
Runtime de agente local & Envuelve manejadores de Python y herramientas locales deterministas como capacidades de agentes administradas por el SDK. \\
Runtime de agente Lingo & Envuelve aplicaciones, flujos, herramientas, máquinas de estado y guardrails opcionales de Lingo como capacidades ordinarias del SDK. \\
Almacén de trazas del SDK & Rastrea eventos de comunicación, fallos en tiempo de ejecución, estado de tareas y observaciones de artefactos. \\
Normalizador de artefactos & Convierte salidas remotas, locales y lingüísticas en registros utilizables por validadores y agregadores. \\
Fábrica de capacidades auxiliares & Usada por \texttt{CoordinationAgent} para especificar agentes lingüísticos temporales y ligeros cuando están justificados. \\
\bottomrule
\end{tabular}
\end{table}

\section{Registro de agentes y frontera de comunicación}

The SDK admits three kinds of agents through one communication boundary. A remote A2A agent is registered with \texttt{register\_a2a\_agent(agent\_card\_url, trust\_level=...)}. The SDK fetches the Agent Card, validates the metadata, extracts identity, service endpoint, skills, modalities, and security requirements, and converts them into an internal registry entry. Because the external representation remains A2A-compatible, other conforming A2A implementations can participate without custom glue.

A local agent is registered with \texttt{register\_local\_agent(name, capabilities, handler, ...)}. The handler may be a Python function, script wrapper, class, symbolic planner, or deterministic tool. The local agent does not need to know about A2A. The SDK supplies the same identity, capability, modality, authority, invocation, and validation metadata that a remote Agent Card would otherwise provide. During dispatch, the SDK invokes the local handler through the Local Agent Runtime and records the result as a normal task observation.

A linguistic agent is registered through the SDK's linguistic-agent registration method. The wrapped object may be a Lingo app, flow, structured generation routine, state machine, tool set, or guardrail-bearing conversational process. The SDK exposes the linguistic agent as an ordinary capability provider. Lingo remains optional and internal to that capability. The \texttt{CoordinationAgent} may invoke the linguistic capability through \texttt{sdk.send\_task(...)} when it needs interpretation, extraction, classification, explanation, or candidate planning.

No registered agent communicates directly with another registered agent through application-level shortcuts. The \texttt{CoordinationAgent} also does not bypass the SDK to reach remote transports, local handlers, or Lingo runtimes. This preserves one communication boundary while keeping planning and feasibility outside the SDK.

\section{Estructuras de datos mínimas}
\label{sec:data-structures}

The implementation can be described through two public surfaces and a set of typed conceptual records. The records are represented as Pydantic models in the prototype so that Lingo structured generation, deterministic validation, SDK registration, coordination-agent planning, and test fixtures use the same schema surface. The records are not tied to a deployment platform; they define the information that the prototype must preserve.

\begin{Verbatim}[fontsize=\small,breaklines=true,breakanywhere=true]
CoordinationSdk:
  register_a2a_agent(agent_card_url, trust_level="standard")
  register_local_agent(name, capabilities, handler, ...)
  register_linguistic_agent(name, capabilities, lingo_app_or_flow, ...)
  registry_snapshot()
  capability_index()
  send_task(agent_id, task)
  trace()
  reset_session()

CoordinationAgent:
  coordinate(user_request, context=None)
  check_feasibility(problem_or_plan)
  trace()

AgentRegistryEntry:
  agent_id
  name
  agent_kind
  service_endpoint
  invocation_endpoint
  description
  skills
  input_modes
  output_modes
  status
  trust_level
  validation_contract
  source_card

CapabilityRequirement:
  name
  description
  input_schema
  output_schema
  input_modes
  output_modes
  constraints
  auxiliary_eligible
  required_trust_level

ProblemRequest:
  user_goal
  requirements
  constraints
  required_artifacts
  context

TaskSpec:
  task_id
  requirement_name
  assigned_to
  auxiliary_spec_id
  depends_on
  expected_artifacts

FeasibilityReport:
  feasible
  matched_agents
  generated_nlp_agents
  missing_capabilities
  risks
  explanation
  evidence

SolutionProposal:
  tasks
  generated_nlp_agents
  selected_agents
  execution_order
  expected_artifacts
  completion_criteria

GeneratedNlpAgentSpec:
  purpose
  input_schema
  output_schema
  method
  validation_rule
  lifecycle
  authority_bounds
  persists

CoordinatorState:
  registry_snapshot
  interpreted_request
  candidate_plan
  feasibility_evidence
  trace

TraceEvent:
  event_type
  message
  data
  timestamp
\end{Verbatim}

The \texttt{CoordinationSdk} public surface is intentionally limited to registration, inspection, dispatch, trace access, and session control. It hides the differences among remote A2A agents, local Python agents, and Lingo-backed linguistic agents while preserving enough metadata for the coordination agent to apply feasibility rules uniformly. The \texttt{CoordinationAgent} public surface contains the complete coordination workflow and feasibility query. The \texttt{AgentRegistry\allowbreak Entry} materializes the discoverable agent set. The \texttt{Capability\allowbreak Requirement} and \texttt{Problem\allowbreak Request} represent the user's goal, constraints, required capabilities, and required artifacts. The \texttt{Solution\allowbreak Proposal} is a candidate plan, not an executable plan until accepted. The \texttt{Feasibility\allowbreak Report} explains whether the coordination agent can solve the problem and why. The \texttt{Generated\allowbreak Nlp\allowbreak Agent\allowbreak Spec} represents a temporary linguistic capability, such as a classifier or extractor, that can be created for the current solution without becoming a permanent system agent. The \texttt{Coordinator\allowbreak State} and \texttt{Trace\allowbreak Event} records provide the auditable ledger required by the traceability property.

\section{Flujo de trabajo del coordinador}

The application-level coordinator is \texttt{CoordinationAgent}. It talks to the distributed system only through \texttt{CoordinationSdk}, but it keeps planning and authorization outside the SDK. This means the SDK is reusable for non-coordinating clients, while the thesis contribution remains concentrated in the coordination agent's feasibility-aware workflow.

\begin{Verbatim}[fontsize=\small,breaklines=true,breakanywhere=true]
sdk = CoordinationSdk(...)
analyzer = FeasibilityAnalyzer(...)
agent = CoordinationAgent(
    sdk=sdk,
    feasibility_analyzer=analyzer,
    linguistic_interpreter="optional-sdk-agent-id",
    auxiliary_factory=optional_factory,
    artifact_aggregator=aggregator,
)

sdk.register_a2a_agent(agent_card_url, trust_level="standard")
sdk.register_local_agent(name, capabilities, handler, ...)
sdk.register_linguistic_agent(name, capabilities, lingo_app_or_flow, ...)

result = agent.coordinate(user_request, context=session_context)
\end{Verbatim}

Inside \texttt{coordinate}, the \texttt{CoordinationAgent} performs the simple-first decision process. It asks the SDK for the current registry and capability index, transforms or admits the user's message into a typed \texttt{ProblemRequest}, attempts direct capability matching, requests candidate decomposition only when direct delegation is insufficient, checks the proposal with the \texttt{FeasibilityAnalyzer}, adds a bounded auxiliary specification only when the feasibility model permits it, and dispatches authorized tasks through the SDK.

\begin{Verbatim}[fontsize=\small,breaklines=true,breakanywhere=true]
CoordinationAgent.coordinate(problem):
  registry = sdk.registry_snapshot()
  capabilities = sdk.capability_index()
  request = interpret_or_admit_problem(problem, sdk)

  candidate_plan = direct_or_decomposed_plan(request, capabilities, sdk)
  report = feasibility_analyzer.check(request, registry, candidate_plan)

  if report.has_auxiliary_eligible_gap:
      auxiliary = auxiliary_factory.specify(report.gap)
      candidate_plan = candidate_plan.with_auxiliary_agent(auxiliary)
      report = feasibility_analyzer.check(request, registry, candidate_plan)

  if report.is_not_feasible:
      return structured_infeasibility(report)

  results = [
      sdk.send_task(task.assigned_to, task)
      for task in candidate_plan.execution_order
  ]
  return artifact_aggregator.combine(results)
\end{Verbatim}

Task dispatch remains polymorphic inside the SDK but not visible to the coordination agent. If the selected executor is a remote A2A agent, the A2A Gateway sends the task and normalizes the remote response. If it is a local agent, the Local Agent Runtime invokes the registered handler and converts the result into the same task observation shape. If it is a Lingo-backed agent, the Lingo Agent Runtime executes the approved app, flow, tool, or state machine and returns a normalized artifact. In all cases, planning infeasibility is structured data returned by the \texttt{CoordinationAgent}. SDK exceptions are reserved for invalid SDK configuration or genuine transport/runtime failures.

This workflow keeps the LLM useful but bounded. Lingo may supply typed context handling and structured proposal mechanisms through SDK-registered linguistic capabilities, but it does not alone decide feasibility. Deterministic checks remain responsible for explicit properties such as availability, modality compatibility, declared skills, authority, and missing capabilities. Lingo's native tool-calling interface may be useful for controlled experiments that intentionally remove symbolic checks, but the main hybrid implementation uses developer-controlled structured output so that symbolic authorization remains explicit.

\section{Especializaciones de agentes locales y lingüísticos}

The SDK treats local and linguistic agents as specializations of the same communication boundary rather than as separate architectural paths. A local agent is appropriate when the capability is deterministic, symbolic, scriptable, or already implemented in Python. Examples include a rule-based planner, a schema validator, a database wrapper, a file transformer, or a mathematical tool. The SDK supplies the communication wrapper, so the local agent author focuses on the capability contract and handler behavior.

A linguistic agent is appropriate when the capability transforms, interprets, classifies, summarizes, or produces language-mediated artifacts. It may be implemented with Lingo, an LLM, a smaller NLP model, rules, retrieval, or a hybrid pipeline. When a registered linguistic capability is implemented with Lingo, it follows the primitive mapping described at the beginning of this chapter; other implementations must still expose the same typed capability, artifact, and validation contracts to the SDK.

Generated NLP agents are solution-local. The \texttt{CoordinationAgent} includes their specifications in the \texttt{SolutionProposal}, including purpose, method, expected input, expected output, validation rule, and lifecycle. They do not automatically enter the registry. Promotion to a persistent registered agent would require a separate SDK admission process.

\section{Diagrama de arquitectura}

Figure~\ref{fig:implementation-architecture} summarizes the implementation structure. The diagram is intentionally abstract enough to support later coding choices while still showing that coordination logic belongs to \texttt{CoordinationAgent}, while communication and runtime adaptation belong to \texttt{CoordinationSdk}.

\begin{figure}[htbp]
\centering
\fbox{
\begin{minipage}{0.90\textwidth}
\centering
Problema del usuario o solicitud de aplicación\\
$\downarrow$\\
\texttt{CoordinationAgent}\\
$\downarrow$\\
\begin{tabular}{ccc}
Interpretación & Viabilidad & Agregación \\
\end{tabular}\\
$\downarrow$\\
Frontera de middleware \texttt{CoordinationSdk}\\
$\downarrow$\\
\begin{tabular}{ccc}
Pasarela A2A & Runtime local & Runtime Lingo \\
\end{tabular}\\
$\downarrow$\\
\begin{tabular}{ccc}
A2A remotos & Python locales & Lingüísticos \\
\end{tabular}\\
$\downarrow$\\
Artefactos, estados, fallos y traza de comunicación\\
$\downarrow$\\
Resultado final o informe estructurado de inviabilidad
\end{minipage}
}
\caption{\texttt{CoordinationAgent} sobre middleware SDK para coordinación compatible con A2A entre agentes remotos, locales y lingüísticos.}
\label{fig:implementation-architecture}
\end{figure}

\section{Secuencia de ejecución}

The execution sequence begins when the user or application submits a problem request to \texttt{CoordinationAgent.coordinate}. The coordination agent asks the SDK for the current registry snapshot and capability index. It obtains a typed problem representation from a structured caller or by invoking a registered linguistic capability through the SDK. The coordination agent then attempts direct matching before considering decomposition. The \texttt{FeasibilityAnalyzer} checks the candidate plan and either approves it, permits an auxiliary NLP specification to be added, or rejects it as infeasible. Approved tasks are dispatched through \texttt{sdk.send\_task(...)}. The SDK chooses the correct internal runtime, records status and failures, normalizes artifacts, and exposes communication trace events. The coordination agent aggregates the authorized outputs or returns an infeasibility report.

\begin{figure}[htbp]
\centering
\fbox{
\begin{minipage}{0.88\textwidth}
\begin{enumerate}
\item El usuario o la aplicación envía un problema y restricciones a \texttt{CoordinationAgent}.
\item El agente de coordinación lee el registro y las capacidades mediante \texttt{CoordinationSdk}.
\item El agente de coordinación obtiene o construye un \texttt{ProblemRequest} tipado.
\item El agente de coordinación intenta correspondencia directa de capacidades antes de la descomposición.
\item El agente de coordinación solicita planificación candidata solo cuando la correspondencia directa es insuficiente.
\item El analizador de viabilidad valida habilidades, modalidades, dependencias, autoridad y contratos de evidencia.
\item La fábrica de capacidades auxiliares añade una especificación temporal solo cuando está justificada.
\item El agente de coordinación despacha tareas autorizadas mediante \texttt{sdk.send\_task(...)}.
\item El SDK registra estado en tiempo de ejecución, errores, artefactos y evidencia de comunicación.
\item El agente de coordinación devuelve el resultado final o un informe estructurado de inviabilidad.
\end{enumerate}
\end{minipage}
}
\caption{Secuencia de ejecución para la resolución de problemas mediada por \texttt{CoordinationAgent}.}
\label{fig:execution-sequence}
\end{figure}

\section{Manejo de fallos e inviabilidad}

The implementation treats infeasibility as a first-class coordination-agent result. If no remote, local, linguistic, or approved auxiliary capability can satisfy a required subtask, the \texttt{CoordinationAgent} returns a structured report. If an agent is unavailable, the coordination agent may select another compatible agent from the SDK registry; if no replacement exists, it reports the missing dependency. If modalities are incompatible, the report identifies which input or output mode cannot be bridged. If a linguistic planner proposes a decomposition that violates explicit constraints, the \texttt{FeasibilityAnalyzer} rejects it.

The SDK normalizes runtime failures rather than planning failures. A2A transport failures, local handler exceptions, Lingo execution failures, invalid artifacts, cancellations, and timeouts are recorded as structured observations linked to the relevant task and agent. This behavior is essential for simplicity and trust. The system should not create complex plans merely to appear capable. It should coordinate only when the available and generated capabilities justify the solution. Otherwise, the coordination agent should explain the boundary clearly. Lingo guardrails may detect obvious unsafe, irrelevant, or ambiguous inputs early, but such reflexes are treated as advisory preprocessing rather than as replacements for the feasibility verdict.

\section{Síntesis del capítulo}

The design specifies a distributed multi-agent environment with a clear split between coordination and communication. \texttt{CoordinationAgent} owns planning, feasibility, authorization, auxiliary decisions, and aggregation. \texttt{CoordinationSdk} owns A2A interoperability, local runtime invocation, Lingo runtime invocation, registry normalization, task dispatch, artifact normalization, communication failures, and trace access. Agents remain independently deployable and vendor independent. A2A provides the interoperability backbone through Agent Cards and task exchange. Local Python agents and Lingo-backed linguistic agents are wrapped by the SDK so that they participate through the same registry, dispatch, artifact, and trace mechanisms as remote A2A agents. The evaluation defined in Chapters~\ref{chap:methodology} and~\ref{chap:results} must determine whether an implementation preserves both communication-boundary compliance and coordination-ownership compliance in practice.

\chapter{Evaluation and Results}
\label{chap:results}

\section{Evidence Status}

This chapter distinguishes corrected version 0.2 preflight evidence from final study evidence. Reports produced by the earlier implementation are legacy artifacts and are excluded from every table and conclusion below. The comparative 36-case local-model study is not yet scored because the independent label sign-off is incomplete. Consequently, this draft reports implementation verification but does not report model accuracy, baseline superiority, or an answer to RQ4 based on the planned comparative study.

\begin{table}[htbp]
\centering
\caption{Corrected version 0.2 preflight gates.}
\label{tab:v02-preflight}
\begin{tabular}{p{0.50\textwidth}rr}
\toprule
\textbf{Gate} & \textbf{Observed} & \textbf{Required} \\
\midrule
Python tests & 78 passing & all passing \\
Overall source coverage & 55\% & at least 85\% \\
Deterministic scenarios & 7/7 & 7/7 \\
Docker A2A checks & 11/11 & 11/11 \\
PostgreSQL distributed checks & 5/5 & 5/5 \\
Official A2A sample interoperability & 1/1 & passing \\
Ruff and mypy & passing & passing \\
Independent corpus sign-off & incomplete & complete \\
Local-model outputs & 0/540 & 540/540 \\
\bottomrule
\end{tabular}
\end{table}

The coverage gate is the principal automated quality blocker. Aggregate coverage is not used to conceal critical gaps: the corrected measurement reports 88\% for the coordination agent, 83\% for feasibility, 77\% for trace validation, 42\% for the PostgreSQL store, and no pytest coverage in the newly added corpus and study runners. Thus the release is not defense-ready even though the executed tests pass.

\section{Fail-Closed Symbolic Boundary}

Regression tests show that a promised artifact does not make a validatorless task verifiable; legacy textual constraints become unresolved typed constraints and force refusal; missing JSON-pointer paths, false predicates, canonical capability-ID mismatch, incomplete mode coverage, and nested JSON Schema type violations are rejected. Effective trust is assigned by local admission policy rather than copied from a remote Agent Card. Cards requiring an unavailable credential scheme are refused, and plain HTTP requires an explicit development exception.

Candidate plans remain non-authoritative. Direct delegation is selected only if it is feasible; otherwise bounded candidates are checked and ranked deterministically by task count, auxiliary count, DAG depth, and stable task identifiers. These tests support prototype correctness of the implemented predicates, not mathematical proof of arbitrary plan correctness.

\section{Dependency, Recovery, and Trace Behavior}

Execution tests confirm that a task is dispatched only after every declared prerequisite completes. Failure, timeout, refusal, or unknown outcome skips descendants while independent branches continue. Consumers receive artifacts only from declared dependencies. On takeover, unsafe unknown work is recorded as \texttt{unknown} and is not blindly redispatched. Completion additionally requires the typed completion contract.

Concurrent-session tests show that returned traces contain only persisted events for their session. Authorization precedes attempt start, SDK transport observations are persisted with session, plan, task, and attempt identity, and a trace validator checks ordering and terminal observations. Registration diagnostics remain outside operational session traces.

\section{Corrected End-to-End Harnesses}

The seven deterministic statuses match their references: three successful completions, one authorized runtime failure, and three pre-dispatch infeasibility decisions. The privacy case now derives email and phone values from transcript text; the answers are no longer supplied as separate fixture fields. The task-local auxiliary lifecycle expires after use.

The corrected Docker harness passes eleven checks covering health, registry discovery, direct and decomposed A2A execution, missing capability, insufficient authority, incompatible modes, invalid artifacts, timeout, registry outage, and session replay. The decomposed case originally exposed a stale assumption that every prior artifact should be forwarded. The corrected planner declares dependencies from canonical mode flow, and the harness now verifies that the final consumer receives only its declared predecessor artifact.

The PostgreSQL harness passes five checks: live lease conflict, takeover after an abandoned read-only attempt, process exit after external dispatch of unknown-side-effect work, stale-fence rejection, and database invariant rejection. It records zero duplicate dispatches and repeated effects in these cases, one lease conflict, one stale-write rejection, terminal correctness, and measured recovery latency. These observations do not establish consensus or cover all crash windows.

An additional test launches the vendored official A2A Hello World sample as an independent server, ingests its Agent Card, preserves its \texttt{echo\_bot} skill identifier, authorizes the matching task, and receives its remote artifact. This establishes one concrete interoperability path; it is not evidence of universal A2A compatibility.

\section{Comparative Study Gate}

The public corpus contains 36 cases and has hash \texttt{b9fd8b6f8398eaf3\allowbreak 3b8676e14288f0d3\allowbreak c4956cca5a9b915c\allowbreak a7b68990167c9c1d}. All three required models are installed and the LM Studio server is reachable. The sign-off record is deliberately marked unapproved. The collection runner therefore refuses to produce the 540 submitted outputs, and the scoring program independently refuses to read hidden labels. This is a successful blinding control, not an empirical result.

RQ1--RQ3 receive implementation-level support from the corrected boundary and harnesses. RQ4 remains open until sign-off, collection, scoring, uncertainty analysis, and a clean evidence freeze are complete. The abstract and conclusion therefore contain no numerical comparative claim.

\chapter{Discussion}
\label{chap:discussion}

\section{Interpretation}

The contribution is an integrative authorization/refusal boundary. Established mechanisms---typed schemas, planning constraints, access policy, task DAGs, leases, fencing, idempotency, and protocol adapters---are combined so that linguistic output cannot silently acquire executable authority. The new contribution claimed here is their operational integration and traceable refinement into one prototype, not invention of each mechanism or formal verification of arbitrary agents.

The LLM is useful where wording varies and decomposition is not known in advance. It is not an authority. Stable capability identifiers, local trust assignment, typed constraints, validation contracts, and completion contracts determine what may be dispatched. A structured rule-only oracle remains a useful upper bound because it shows what can be achieved when interpretation is already supplied; it is not a fair replacement for an end-user system receiving natural language.

Version 0.3 exposed the initial integration's central weakness: safety was obtained through universal refusal. Version 0.4 narrowed the linguistic role and produced substantially better but model-sensitive development results. Version 0.5 then showed that the historical identifier-copying bridge did not transfer to fresh matched pairs: it retained zero unsafe acceptance only by accepting five of 360 feasible repeated observations. That negative experiment motivated the production redesign and is no longer treated as the primary comparison.

The frozen v0.7 study evaluates raw language through production semantic admission and the shared compiler. At primary seed 11, the hybrid attains 71.88\% balanced accuracy and 44.79\% feasible recall, with one unsafe acceptance in 192 observations. The direct LLM attains 53.13\% balanced accuracy and 98.96\% feasible recall but makes 89 unsafe acceptances. The pair-cluster interval for the hybrid-minus-direct unsafe-acceptance rate is entirely negative, so the safety-superiority gate passes. The recall difference is between $-63.54$ and $-45.83$ percentage points, so non-inferiority fails. Safety and utility remain in tension rather than being jointly achieved.

Model separation is essential. Qwen3-1.7B hybrid recall is 56.25\% with no unsafe acceptance; Gemma-4-E2B refuses every feasible case and therefore has zero recall; Qwen3-8B reaches 78.13\% recall but produces the sole unsafe hybrid acceptance. High seed agreement---98.44\% for the hybrid and 96.88\% for direct---shows that the primary pattern is not explained by widespread seed instability. It does not establish backend seed compliance or generalize beyond these quantizations and this laptop runtime.

The exact-name/alias baseline solves all 64 designed cases. Its perfect balanced accuracy means the lexical-superiority gate fails decisively: the hybrid-minus-alias interval is $[-32.81,-23.44]$ percentage points. This is both a negative result for the proposal on the corpus and a construct warning. The designed language contains enough exact aliases for a deterministic matcher to recover the author labels; therefore the corpus does not establish that an LLM is necessary. Future evaluation must include independently sourced utterances whose semantics cannot be recovered through catalog phrase overlap while keeping the same safety boundary.

The production redesign nevertheless isolates where failure occurs. The oracle/compiler invariants all pass, provider selection is registry-order independent, alternative recovery succeeds, and search beyond 4,096 assignments refuses reproducibly. Eight observed primary-seed hybrid acceptances replay through the fixture A2A execution boundary with dependency artifacts and validation contracts, all completing. These results support the implemented deterministic boundary and execution path, not comparative superiority of semantic admission.

\section{Guarantees and Non-Guarantees}

Within the tested implementation, authorization precedes dispatch, unresolved constraints fail closed, declared failed dependencies block descendants, independent branches may continue, session traces are isolated, and stale fenced writes are rejected. These are implementation properties supported by tests and measured traces. They are not mathematical proofs over every deployment, scheduler, transport, or database failure.

Contract validity is weaker than semantic truth. JSON Schema can establish nested type and required-field conformance; a validator can establish a domain-specific property only if that property is encoded. Neither implies that an open-ended summary is complete or that a physical action is safe. The trusted registry and honest-agent assumptions remain explicit.

The admission boundary controls allowed origins, HTTPS policy, locally assigned trust, required credential schemes, and request-header propagation. It does not provide cryptographic identity, attestation, or defense against a malicious admitted agent. Plain HTTP is enabled only for explicit development fixtures.

The legacy PostgreSQL lease-and-fencing mode is not consensus because one database remains its serialization point. The new etcd mode is consensus-backed: Raft orders registry, ownership, plan, attempt, and terminal mutations while a voter majority is available. This does not imply Byzantine tolerance or exactly-once external effects. Fencing prevents an obsolete coordinator from committing authoritative state only when every mutation checks the fence, and external safety additionally requires the receiving agent to enforce the highest fence and durably deduplicate operation keys.

\section{Critical Assessment of the Distributed Implementation}

The strongest design decision is the separation between logical authority and physical availability. The theory still has one accountable authorization role, while etcd removes the single volatile coordinator and single PostgreSQL coordination database from the authoritative path. Transactional owner comparisons, monotonic revision fences, attempt-before-result checks, terminal uniqueness, lease-backed agent registration, and fail-closed quorum errors reinforce one another instead of treating leader election alone as sufficient. Learner-first admission and serialized membership changes reduce reconfiguration risk. Keeping JSONL and PostgreSQL adapters also permits contract comparison and incremental migration rather than an irreversible cutover.

The unified container is operationally simpler, but it creates correlated resource and failure domains. A CPU-intensive planner, slow disk, container restart, host loss, or network namespace fault can affect both the coordinator and its local etcd member. The dispatch concurrency limit mitigates but does not eliminate heartbeat starvation. A production topology must place voters on independent hosts or failure domains, reserve CPU and durable storage for etcd, and monitor disk latency. Three containers on one Docker host demonstrate process-level behavior, not host-level fault tolerance.

Automatic scaling is also easy to overstate. The implementation automatically discovers an existing cluster and reconciles the membership of coordinator nodes that already exist. It does not create compute instances, repair a permanently lost majority, or safely choose a new cluster after loss of quorum. DNS/direct seeds remain necessary across routed networks; multicast is a link-local fallback and may be filtered by container, cloud, or enterprise networks. Target one has no fault tolerance. Targets five and seven tolerate more failures but add write latency and have not yet been exercised by the accepted harness.

The security boundary is appropriate for a trusted private prototype but weak for hostile deployment. HMAC authenticates possession of a cluster-wide secret, not an individual node identity; compromise of one coordinator therefore permits forged cluster messages until the shared secret is rotated. Nonce and timestamp checks depend on bounded clock skew and local replay caches. Discovery authentication does not encrypt or authenticate etcd peer and client transport. Deferring mTLS, certificate rotation, per-node authorization, and attestation leaves etcd ports and shared secrets as high-value assets. The separate agent-registration secret is useful compartmentalization, but admitted malicious agents and compromised voters remain out of scope.

The custom HTTP-gateway client keeps dependencies small and supports endpoint rotation, transactions, leases, and watch resumption, but it also places subtle retry, compaction, cancellation, and error-classification behavior in project code. Quorum-unavailable errors must never be converted into permissive retries or stale reads. The asynchronous PostgreSQL projector correctly avoids a readiness dependency, yet event deletion must remain disabled until projection or export is acknowledged and a retention policy is explicit.

Consensus-v4 shows a materially stronger but still bounded distributed result. Formation at three, five, and seven voters; leader loss; minority partition; majority loss and restoration; concurrent ownership; audit failure; failed-voter replacement; all three crash windows; and the supplementary combined-fault scenario passed in every completed repetition. No executed primary safety invariant failed. Nevertheless, only 140 of 150 primary checks executed because one repetition of each live-reconfiguration path timed out before its five checks ran. These are infrastructure-indeterminate outcomes, not demonstrated safety violations. The campaign also records a dirty source tree and is therefore ineligible as accepted clean evidence. RQ5 is partially supported for the repeatedly completed conditions, while complete reconfiguration reliability and clean-source reproducibility remain open.

\section{Claim and Terminology Discipline}

The writing must distinguish three concepts that earlier drafts conflated. \emph{Centralized authorization} describes one logical decision role; it does not necessarily mean one process. \emph{Consensus-backed replication} describes how that role's state is ordered and recovered; it does not make authorization a peer negotiation. \emph{Automatic membership reconciliation} describes adding or removing already running nodes; it is not automatic compute provisioning. Accordingly, ``fault tolerant'' should always be qualified as crash-fault tolerant while a majority survives, and ``no single point of failure'' should be limited to the authoritative coordination path under independent voter failure domains. The implementation claims duplicate suppression for measured cases, not exactly-once effects.

\section{Threats to Validity}

Construct validity is protected by separating semantic fields, feasibility, execution, artifacts, trace, and recovery measures, but the author-designed and author-labelled corpus creates construct circularity and researcher-bias risk. Freezing labels before confirmatory collection and hiding them during inference reduces leakage, not bias or lack of independent adjudication. Internal validity remains threatened by model nondeterminism, unverifiable backend seed compliance, prompt sensitivity, baseline choices, fixture ownership, one separately classified model-runtime failure, and distributed checks authored with the implementation. External validity is limited to 64 held-out designed cases, three small local models, controlled fixture agents, selected PostgreSQL crash faults, and single-host Docker campaigns. Conclusion validity uses 32 matched pairs, model/category reporting, one primary seed, one stability replication, and pair-clustered bootstrap intervals. Repeated measurements cannot compensate for the small designed corpus, author labels, correlated requested-seed outputs, infrastructure-failed trials, or the absence of independent failure domains.

\section{Defense Readiness}

\enlargethispage{2\baselineskip}
The repository is technically evidence-complete for the v0.7 rebuild: the complete test suite, Ruff, mypy, strengthened production branch-coverage gates, deterministic compiler benchmark, eight fixture-A2A replays, and the full three-model matrix pass their integrity checks. The 45-trial consensus-v4 matrix is complete but not accepted clean evidence because its provenance records the preserved dirty worktree. Defense readiness depends on presenting mixed results without dilution: RQ4 passes only safety superiority and symbolic invariants, while RQ5 is partially supported with two reconfiguration repetitions indeterminate. Independent corpus validation remains a disclosed limitation. Administrative tutor, tribunal, clean-source rerun, and packaging metadata remain release blockers. Independent label review and a broader low-overlap raw-language corpus are required before claiming a neutral benchmark, a superior planning pipeline, or a complete map of supported fault conditions.

Version 0.8 addresses the specific v0.7 semantic failure modes through evidence grounding, deterministic ambiguity, quotation and negation boundaries, non-authoritative retrieval, and constraint-directed provider search. Its complete development matrix passes the qualification gates for both permitted models, but this is intentionally a stopping point rather than a success claim: the same cases informed implementation revision. Qwen3-8B is excluded only from v0.8 and its v0.7 result remains unchanged. A cross-version comparison will be descriptive even after confirmation because the corpus and semantic interface changed. Until the two human labeling passes are completed and the 768 held-out observations are validated, v0.8 cannot alter the thesis's confirmatory conclusions.

\chapter{Conclusions and Future Work}
\label{chap:conclusion}

This thesis addressed the gap between linguistically plausible coordination and justified executable coordination in heterogeneous autonomous systems. With respect to the first objective, the state-of-the-art analysis showed that classical multi-agent systems, symbolic AI, LLM agents, and interoperability protocols contribute necessary but individually insufficient mechanisms. Their unresolved intersection is the explicit authorization boundary between language-mediated proposals and coordinated action.

With respect to the second objective, the thesis defined a protocol-independent framework in which capabilities, artifacts, dependencies, authority, commitments, validators, and completion evidence are explicit. Its central result is a total decision at the planning boundary: the coordinator produces either an authorized plan supported by predicate evidence or a structured infeasibility report. With respect to the third objective, that framework was mapped to a Python \texttt{CoordinationAgent} design over \texttt{CoordinationSdk} middleware, with protocol adaptation, registry admission, capability indexing, local and linguistic agent runtimes, feasibility analysis, task monitoring, and artifact aggregation handled at the implementation layer. The mapping illustrates implementability without making any particular communication protocol, linguistic runtime, or local runtime part of the theoretical claim.

With respect to the fourth objective, the thesis executed a reproducible prototype evaluation based on feasible and infeasible deterministic scenarios, Docker A2A system checks, local LLM reference batches, and separate outcome, process, communication-boundary, coordination-ownership, and trace metrics. The evidence shows that the implementation can authorize feasible plans, refuse missing or unauthorized requirements before dispatch, separate runtime failure from planning infeasibility, preserve SDK-mediated communication, and record model-specific linguistic reference behavior. The thesis does not claim production readiness, guaranteed runtime success, or broad empirical superiority over all possible baselines. This boundary is itself consistent with the proposed framework: conclusions should be authorized by available evidence rather than by plausible expectation.

The principal contribution is therefore a precise integration pattern for controlled openness: language supports interpretation and proposal; symbols support authorization and traceability; independently deployed agents retain execution autonomy; and unsupported requests are refused explicitly. The framework answers RQ1--RQ3 at the level of design, formal specification, and prototype behavior. RQ4 is answered for the implemented evidence set by the reported outcome, trace, boundary, latency, and local-model metrics, while larger baseline comparisons remain future work.

Future work should extend the implemented \texttt{CoordinationAgent}, \texttt{CoordinationSdk}, deterministic scenario runner, JSONL ledger, Docker A2A harness, and local LLM reference runner into a larger paired baseline and ablation study. Subsequent research should evaluate dynamic registry changes, adversarial capability advertisements, cryptographic identity and attestation, replicated coordinators with lease-based hot standby, domain-specific semantic validators, human authorization for high-impact actions, and third-party A2A implementations. Evaluation across protocols, local non-linguistic agents, Lingo-authored linguistic agents, larger local and hosted language models, and real operational settings is necessary to determine which findings generalize beyond the first agent-over-SDK instantiation.


\backmatter
\bibliographystyle{plain}
\bibliography{references/references}

\end{document}
